%% 46장 미분 없는 최적화
\chapter{미분 없는 최적화}
\label{ch:derivative-free}

\abstract{기울기를 계산할 수 없거나 믿을 수 없는 함수도 최소화해야 할 때가 있다. 시뮬레이션의 결과나 반올림 잡음이 섞인 값, 꺾인 곳이 있는 함수가 그렇다. 이 장에서는 그런 문제에 가장 널리 쓰이는 Nelder--Mead 방법을 다룬다. 단체를 반사, 확장, 수축, 축소의 네 연산으로 굴려 가는 이 방법은 단순하고 대개 잘 듣지만, 이론적 보장은 놀랄 만큼 약하다. 2차원의 강볼록 함수에서도 정류점이 아닌 점으로 수렴하는 McKinnon의 반례를 그 기하로 설명하고, 반례의 점들이 정확히 $\lambda^k$의 꼴로 줄어드는 것을 수치로 확인한다. 차원이 커지면 표준 계수의 Nelder--Mead가 무너지고 Gao와 Han의 적응 계수가 그것을 크게 늦춘다는 것도 확인한다. 끝으로 차분 기울기와 패턴 탐색처럼 수렴을 보장하는 대안을 정리한다.}

\section{도함수가 없을 때}
\label{sec:no-derivatives}

44장과 45장의 방법은 모두 기울기를 쓴다. 기울기가 없으면 가장 먼저 떠오르는 대안은 33장의 유한 차분이다. 앞쪽 차분으로 기울기 하나에 함수값 $n$번을 더 쓰고 간격을 $h\approx\sqrt u\,\max(1,|x_i|)$로 잡으면 기울기의 상대 정확도가 대략 $\sqrt u\approx10^{-8}$이므로, 준Newton 방법을 그대로 쓸 수 있다. 실제로 scipy의 BFGS에 기울기를 주지 않으면 이렇게 동작하며, 2차원 Rosenbrock 함수를 $(-1.2,1)$에서 $32$번 반복, 함수값 $117$번으로 풀어 오차가 $9.1\times10^{-6}$이다.

차분이 통하지 않는 경우가 있다. 함수값에 반올림이나 몬테카를로의 잡음이 섞여 있으면 차분은 잡음의 차이를 $h$로 나눈 것이어서 쓸모없어진다. 함수가 꺾이는 곳이 있으면 거기서 기울기가 정의되지 않는다. 함수값 하나가 몇 분씩 걸리는 시뮬레이션이라면 기울기 하나에 $n$번을 더 쓰는 것도 부담이다. 이런 문제를 위해 함수값의 대소만 비교하는 \term{직접 탐색}{direct search} 방법들이 있다. 43장의 황금분할이 한 변수의 직접 탐색이었다. 여러 변수로는 Nelder--Mead 방법이 가장 널리 쓰인다.

\section{Nelder--Mead 방법}
\label{sec:nelder-mead}

\subsection{단체의 네 연산}

$\mathbb R^n$의 점 $n+1$개 $x_0,\dots,x_n$이 한 초평면에 있지 않으면 그 볼록 껍질을 \term{단체}{simplex}라 한다. 2차원에서는 삼각형, 3차원에서는 사면체이다. Nelder--Mead 방법은 반복마다 꼭짓점을 함수값 순서 $f(x_0)\le f(x_1)\le\dots\le f(x_n)$로 정렬하고, 가장 나쁜 꼭짓점 $x_n$을 나머지의 무게중심 $\bar x=\sum_{i<n}x_i/n$ 반대편으로 옮겨 보는 일을 되풀이한다. 옮기는 점은 모두 직선 $\bar x+t(\bar x-x_n)$ 위에 있다.
\begin{itemize}
\item \term{반사}{reflection}: $x_r=\bar x+\rho(\bar x-x_n)$, 기본 $\rho=1$.
\item \term{확장}{expansion}: $x_e=\bar x+\rho\chi(\bar x-x_n)$, 기본 $\chi=2$.
\item \term{바깥 수축}{outside contraction}: $x_c=\bar x+\rho\psi(\bar x-x_n)$, 기본 $\psi=1/2$.
\item \term{안쪽 수축}{inside contraction}: $x_{cc}=\bar x-\psi(\bar x-x_n)$.
\item \term{축소}{shrink}: 가장 좋은 꼭짓점 $x_0$을 남기고 나머지를 $x_i\gets x_0+\sigma(x_i-x_0)$로 당긴다. 기본 $\sigma=1/2$.
\end{itemize}
어느 연산을 할지는 반사점의 함수값 $f_r=f(x_r)$이 정한다(\cref{alg:nelder-mead}). 반사점이 가장 좋은 꼭짓점보다도 좋으면 그 방향으로 더 멀리 확장해 보고, 둘 가운데 좋은 쪽을 받아들인다. 가장 좋은 것과 두 번째로 나쁜 것 사이이면 반사점을 그대로 받아들인다. 두 번째로 나쁜 것보다 나쁘면 수축하는데, 반사점이 가장 나쁜 꼭짓점보다는 좋으면 바깥으로, 그보다도 나쁘면 안쪽으로 수축한다. 수축점도 실패하면 단체 전체를 가장 좋은 꼭짓점 쪽으로 축소한다. 그래서 단체는 내리막으로는 늘어나며 굴러가고, 골짜기에서는 그 방향으로 길쭉해지며, 최소점 근처에서는 오그라든다.

\begin{algorithm}[t]
\begin{algorithmic}[1]
\Require 함수 $f$, 처음 단체 $x_0,\dots,x_n$, 계수 $\rho,\chi,\psi,\sigma$, 허용오차 $\mathrm{xatol}$, $\mathrm{fatol}$
\Ensure 근사 최소점 $x_0$
\Loop
  \State $f(x_0)\le\dots\le f(x_n)$이 되게 정렬한다
  \If{$\max_i\|x_i-x_0\|_\infty\le\mathrm{xatol}$이고 $\max_i|f(x_i)-f(x_0)|\le\mathrm{fatol}$} \Return $x_0$ \EndIf
  \State $\bar x\gets\sum_{i<n}x_i/n$, $x_r\gets\bar x+\rho(\bar x-x_n)$
  \If{$f(x_r)<f(x_0)$}
    \State $x_e\gets\bar x+\rho\chi(\bar x-x_n)$, $x_n\gets$ ($f(x_e)<f(x_r)$이면 $x_e$, 아니면 $x_r$)
  \ElsIf{$f(x_r)<f(x_{n-1})$}
    \State $x_n\gets x_r$
  \ElsIf{$f(x_r)<f(x_n)$}
    \State $x_c\gets\bar x+\rho\psi(\bar x-x_n)$; $f(x_c)\le f(x_r)$이면 $x_n\gets x_c$, 아니면 축소
  \Else
    \State $x_{cc}\gets\bar x-\psi(\bar x-x_n)$; $f(x_{cc})<f(x_n)$이면 $x_n\gets x_{cc}$, 아니면 축소
  \EndIf
  \State 축소: $i=1,\dots,n$에 대해 $x_i\gets x_0+\sigma(x_i-x_0)$
\EndLoop
\end{algorithmic}
\caption{Nelder--Mead 방법(scipy의 순서)}
\label{alg:nelder-mead}
\end{algorithm}

한 반복의 비용은 함수값 한두 번이고, 축소할 때만 $n$번이 더 든다. 처음 단체는 시작점 $x$와, 그 성분 하나씩을 $5\%$ 늘린 점 $n$개로 만든다(성분이 $0$이면 $0.00025$로 둔다). 멈춤 판정은 꼭짓점들이 위치로도 함수값으로도 가장 좋은 꼭짓점에서 $10^{-4}$ 이내로 모였을 때이고, 함수값을 $200n$번 쓰면 멈춘다. 2차원 Rosenbrock 함수를 $(-1.2,1)$에서 시작하면 이 기본값으로 $85$번 반복, 함수값 $159$번에 오차 $4.2\times10^{-5}$로 멈춘다. 허용오차를 $\mathrm{xatol}=10^{-8}$, $\mathrm{fatol}=10^{-12}$로 좁히면 $117$번 반복, $219$번으로 오차 $1.6\times10^{-9}$이다. 기울기를 쓰는 BFGS의 $39$번보다는 많지만 차분 BFGS의 $117$번과는 비슷한 수준이다.

\begin{aside}[단체 방법의 두 갈래]
Spendley, Hext, Himsworth는 1962년에 공장의 공정 조건을 실험으로 개선하는 방법으로 정단체를 반사해 나가는 방법을 내놓았다\cite{spendley1962}. 단체의 모양은 언제나 정단체이고 크기만 가끔 줄였다. Nelder와 Mead는 1965년에 확장과 수축을 더해 단체가 함수의 모양에 맞추어 늘어나고 줄어들게 했다\cite{nelder1965}. 이 방법은 MATLAB의 \code{fminsearch}, R의 \code{optim}의 기본값으로 들어가 있다. 선형계획법의 심플렉스 방법(48장)과는 이름만 같고 아무 관계가 없다.
\end{aside}

\section{수렴에 대해 알려진 것}
\label{sec:nm-theory}

Nelder--Mead는 50년 넘게 쓰였지만 그 수렴에 대해 증명된 것은 많지 않다. 걸음이 함수값의 대소만으로 정해지고 충분 감소 같은 조건이 없으므로, 44장의 Zoutendijk 정리나 45장의 Cauchy 감소 같은 논법이 통하지 않는다. Lagarias, Reeds, M.\ Wright, P.\ Wright는 1998년에 표준 계수의 Nelder--Mead를 강볼록 함수에서 분석했다\cite{lagarias1998}. 1차원에서는 최소점으로 수렴함을 보였지만, 2차원에서는 꼭짓점의 함수값들이 같은 값으로 수렴하고 단체의 지름이 $0$으로 간다는 데까지만 보였다. 그 극한이 최소점이라는 것은 증명하지 못했다. 같은 해에 McKinnon이 그 이유를 보여 주었다.

\subsection{McKinnon의 반례}

McKinnon은 2차원의 강볼록이고 연속 미분 가능한 함수에서 Nelder--Mead가 정류점이 아닌 점으로 수렴하는 예를 만들었다\cite{mckinnon1998}. 함수는
\[
f(x,y)=\begin{cases}\theta\phi|x|^\tau+y+y^2,&x\le0,\\\theta x^\tau+y+y^2,&x>0\end{cases}
\]
이고, 여기서는 $\tau=2$, $\theta=6$, $\phi=60$을 쓴다. $x<0$ 쪽이 $60$배 가파르지만 $x=0$에서 연속 미분 가능하고, 최소점은 $(0,-1/2)$이며 최솟값은 $-1/4$이다. 처음 단체를
\[
v_0=(0,0),\qquad v_1=(\lambda_1,\lambda_2),\qquad v_2=(1,1),\qquad\lambda_{1,2}=\frac{1\pm\sqrt{33}}8
\]
로 잡는다. $\lambda_1=0.8431$, $\lambda_2=-0.5931$이고 둘은 $4\lambda^2-\lambda-2=0$의 두 근이다.

이 단체에서 일어나는 일은 다음과 같다. 꼭짓점을 $v_0=(0,0)$, $a_k=(\lambda_1^{k+1},\lambda_2^{k+1})$, $b_k=(\lambda_1^k,\lambda_2^k)$로 쓰자. $k=0$이 처음 단체이다. 가장 나쁜 꼭짓점은 $b_k$이고 무게중심은 $\bar x=a_k/2$이다. 반사점 $\bar x+(\bar x-b_k)=a_k-b_k$의 $x$ 성분은 $\lambda_1^k(\lambda_1-1)<0$이어서 가파른 쪽에 떨어지고, McKinnon의 매개변수에서는 그 값이 $f(b_k)$보다 커서 안쪽 수축이 일어난다. 안쪽 수축점은 성분마다
\[
\frac12\cdot\frac{\lambda^{k+1}}2+\frac12\lambda^k=\lambda^k\Bigl(\frac\lambda4+\frac12\Bigr)=\lambda^{k+2}
\]
이다. 마지막 등호가 $4\lambda^2=\lambda+2$, 곧 $\lambda_1$과 $\lambda_2$를 정한 방정식이다. 그래서 수축점은 정확히 $(\lambda_1^{k+2},\lambda_2^{k+2})$이고, 그것이 $b_k$보다 좋아 받아들여지면 새 단체는 $v_0$, $a_{k+1}$, $b_{k+1}$로 처음과 같은 꼴이다. McKinnon은 함수값의 대소가 매번 이렇게 정해지는 조건을 매개변수에 대해 구했다. 이 반복은 끝없이 이어지고 두 꼭짓점이 $v_0=(0,0)$으로 다가가므로 단체 전체가 $(0,0)$으로 오그라든다. 그런데 $(0,0)$에서 기울기는 $(0,1)$이다. 가장 좋은 꼭짓점 $v_0$은 한 번도 바뀌지 않는다.

scipy 1.18.1로 이 단체에서 시작하면 정확히 그렇게 된다. 함수값을 부른 점을 차례로 보면 처음 세 꼭짓점 뒤로 반사점 $(-0.157,-1.593)$, 수축점 $(0.7108,0.3517)$, 반사점 $(-0.132,0.945)$, 수축점 $(0.5992,-0.2086)$이 이어지는데, 두 수축점은 $(\lambda_1^2,\lambda_2^2)$, $(\lambda_1^3,\lambda_2^3)$과 반올림 오차 안에서 같다. 허용오차 $10^{-8}$에서 $109$번 반복, 함수값 $219=3+2\times108$번으로 멈추어 모든 반복이 반사 하나와 수축 하나였음을 보여 준다. 돌려준 점은 $(0,0)$이고 함수값은 $0$으로, 참 최솟값 $-1/4$에 이르지 못했다. 단체가 기울기와 거의 직각인 방향으로 납작해지면서 오그라든 것이다. $\lambda_2<0$이므로 $y$ 성분은 부호를 바꾸며 줄고, 내리막 방향 $-y$로는 한 번도 나아가지 못한다.

이 반례는 Nelder--Mead의 약점을 그대로 드러낸다. 단체가 납작해지면 꼭짓점들이 거의 한 직선 위에 있어 무게중심 반대편으로의 반사가 내리막 방향을 탐색하지 못한다. 충분 감소 조건이 없으므로 이런 퇴화를 막는 장치가 없다. Kelley는 단체의 모양으로 만든 단체 기울기가 충분 감소를 주지 못하면 정체로 보고 단체를 다시 만드는 방법을 제안했다\cite{kelley1999}. 실용적으로는 멈춘 점에서 Nelder--Mead를 한 번 더 시작하는 것이 같은 효과를 낸다.

\section{차원이 커질 때}
\label{sec:nm-dimension}

Nelder--Mead는 차원이 커질수록 급격히 나빠진다. 이차 함수 $f(x)=\sum_{i=1}^ni\,x_i^2$를 $x=(1,\dots,1)$에서 시작해 허용오차를 $10^{-12}$로 두고 돌린 결과가 \cref{tab:nm-dimension}이다. 표준 계수로는 $n=20$에서 이미 함수값이 $44{,}222$번 들고, $n=40$에서는 $10^6$번을 써도 $f=5.85$에 머문다. 단체가 $n+1$개의 꼭짓점으로 $n$차원의 모양을 따라가야 하는데, 반사와 수축이 한 번에 꼭짓점 하나만 움직이므로 단체가 함수의 모양에 맞추어지는 데 시간이 오래 걸리고, 그 사이에 단체가 납작해지기 쉽다.

Gao와 Han은 계수를 차원에 맞추어
\[
\rho=1,\qquad\chi=1+\frac2n,\qquad\psi=\frac34-\frac1{2n},\qquad\sigma=1-\frac1n
\]
로 두는 \term{적응 계수}{adaptive parameters}를 제안했다\cite{gao2012}. $n=2$이면 표준 계수와 같다. 차원이 크면 확장이 덜 과격하고 수축과 축소가 덜 급해져서 단체가 납작해지는 것을 늦춘다. \cref{tab:nm-dimension}에서 적응 계수는 $n=40$도 함수값 $24{,}034$번으로 푼다. 다만 $n=5$에서는 표준 계수가 조금 낫듯이 언제나 이기는 것은 아니다.

\begin{table}[t]
\begin{sdtable}{colspec={lccccc}, row{1}={font=\small\bfseries}}
\toprule
차원 $n$ & $2$ & $5$ & $10$ & $20$ & $40$ \\
\midrule
표준 계수 & $198$ & $705$ & $3{,}283$ & $44{,}222$ & $10^6$(실패) \\
적응 계수 & $198$ & $877$ & $2{,}384$ & $6{,}497$ & $24{,}034$ \\
\bottomrule
\end{sdtable}
\caption{$f(x)=\sum_ii\,x_i^2$를 $(1,\dots,1)$에서 시작해 $\mathrm{xatol}=\mathrm{fatol}=10^{-12}$까지 Nelder--Mead로 최소화할 때 함수값을 계산한 횟수(scipy 1.18.1). 표준 계수의 $n=40$은 $10^6$번 뒤에도 $f=5.85$였다.}
\label{tab:nm-dimension}
\end{table}

\section{수렴을 보장하는 대안}
\label{sec:pattern-search}

수렴의 보장이 필요하면 다른 직접 탐색을 쓴다. \term{패턴 탐색}{pattern search}은 현재 점에서 정해진 방향 집합(예컨대 $\pm e_i$)으로 걸음 $\Delta$만큼 가 본 점들을 비교해, 나아지는 점이 있으면 옮기고 없으면 $\Delta$를 줄인다. 방향 집합이 공간 전체를 양의 일차 결합으로 덮으면(예컨대 $\pm e_i$의 $2n$개), 어느 점에서든 그 가운데 하나는 기울기와 예각을 이룬다. 그래서 $f$가 연속 미분 가능하면 걸음이 실패할 때마다 기울기가 $\Delta$에 비례하는 크기로 작다는 것을 알 수 있고, $\Delta\to0$인 부분열에서 기울기가 $0$으로 간다. Torczon은 1997년에 이런 방법들을 한 틀로 묶어 이 수렴을 증명했다\cite{torczon1997}. 단체의 모양이 퇴화하지 않도록 방향 집합을 고정하는 것이 Nelder--Mead와의 차이이다. 그 대가로 대개 더 많은 함수값을 쓴다.

또 하나의 길은 함수값으로 이차 모형을 보간해 45장의 신뢰 영역 방법에 넣는 \term{모형 기반 방법}{model-based method}이다. Powell의 NEWUOA\cite{powell2006newuoa}와 그 경계 있는 판 BOBYQA가 대표적이며, 매끄러운 함수에서는 Nelder--Mead보다 훨씬 적은 함수값으로 수렴한다. 함수가 매끄럽고 잡음이 없다면 앞에서 본 차분 BFGS도 좋은 선택이다. Nelder--Mead가 여전히 쓰이는 까닭은 구현이 짧고, 꺾인 함수와 잡음에 둔하며, 낮은 차원에서는 대개 잘 듣기 때문이다.

\begin{instardust}[\code{docs/design/optimize.md} §3.5]
\code{sd.optimize.minimize(method='nelder-mead')}는 scipy의 \code{\_minimize\_neldermead}를 옮긴 것이다(\code{csrc/optimize/nelder\_mead.cpp}). 처음 단체(성분마다 $1.05$배, $0$이면 $0.00025$)나 사용자의 \code{initial\_simplex}, \code{adaptive}의 Gao--Han 계수, 멈춤 판정, 함수값 상한을 넘는 호출을 하지 않고 그 반복을 멈추는 것까지 같다. 경계가 있으면 위 경계를 넘은 처음 꼭짓점을 경계에 대해 반사한 뒤 자르고, 새 점은 만들 때마다 상자 안으로 자른다. 결과를 비트까지 맞추려면 산술뿐 아니라 정렬도 같아야 한다. numpy의 \code{argsort}는 기본이 introsort라 원소가 $17$개 미만이면 삽입 정렬이어서 같은 값의 순서가 유지되지만 그 이상이면 유지되지 않으므로, 함수값이 같은 꼭짓점들의 순서까지 맞추려고 그 알고리즘을 그대로 옮겼다. 무게중심은 numpy처럼 $+0.0$에서 행을 차례로 더한다. 정렬은 행을 옮기지 않고 행 번호만 옮긴다.
\end{instardust}

\section{맺음}

미분 없는 최적화는 정보가 가장 적은 상황에서의 최적화이다. 함수값의 대소만 쓰는 방법은 기울기를 쓰는 방법보다 느리고, Nelder--Mead처럼 실용적으로 성공한 방법도 이론적인 보장은 약하다. McKinnon의 반례는 단체가 납작해지면 내리막을 보지 못한다는 기하적 약점을 정확히 짚었고, 그 약점은 차원이 커질 때 더 두드러진다. 보장이 필요하면 방향 집합을 고정한 패턴 탐색이나 모형 기반 신뢰 영역을 쓰고, 함수가 매끄러우면 차분 기울기로 44장의 방법을 쓰는 것이 낫다. 다음 장에서는 다시 도함수를 쓰되, 목적 함수가 잔차의 제곱합이라는 구조를 이용하는 비선형 최소제곱으로 넘어간다.

직접 탐색의 이론과 실제는 Conn, Scheinberg, Vicente의 책\cite{conn2009dfo}과 Kolda, Lewis, Torczon의 개관\cite{kolda2003}이 잘 정리했고, Nelder--Mead의 수렴 문제는 Lagarias 등\cite{lagarias1998}과 McKinnon\cite{mckinnon1998}의 논문에 있다.

\chapterreferences
